\begin{center}
    1. ТЕКСТ ЗАДАНИЯ
\end{center}

Цель работы: дальнейшее изучение процессора с регистровой архитектурой, изучение форматов команд и данных, ознакомление
с работой программы-симулятора. Выполнение загрузки команд и данных, 
выполнение простейших программ арифметико-логической обработки данных из памяти, использование косвенной адресации.

Практическая часть работы включает выполнение следующих действий:
\begin{enumerate}
    \item формирование адресов числовых значений в соответствии с индивидуальным заданием, перевод их в шестнадцатеричную систему счисления,
    \item запись целочисленных данных в память по заданным адресам,
    \item по заданному алгоритму составление и выполнение простейшей программы работы с целочисленными данными, хранящимися в памяти, 
        с использованием различных способов косвенной адресации,
    \item выполнение программы,
    \item контроль результатов работы программы.
\end{enumerate}

Правильность разработки и выполнения программ арифметико-логической обработки данных 
контролируется путем ручной трассировки заданных алгоритмов 
с последующим сравнением результатов работы программ с результатами ручной трассировки